% =====================================================================
%  Jonathan Woollett-Light -- Curriculum Vitae
%
%  This file is the SOURCE OF TRUTH for the CV content.
%
%  It is real LaTeX: `pdflatex cv.tex` produces a typeset PDF.
%  The website (index.html) is GENERATED from this file by pandoc,
%  using build/template.html for the page chrome (theme toggle, print
%  button) and the stylesheet.  Run `npm run build` after editing.
%  Do not edit index.html by hand; your changes will be overwritten.
% =====================================================================
\documentclass[11pt,a4paper]{article}

\usepackage[margin=2.2cm]{geometry}
\usepackage[hidelinks]{hyperref}

\title{Jonathan Woollett-Light}
\author{}
\date{}

% ---------------------------------------------------------------------
%  Semantic macros.
%
%  pandoc expands these when generating the website, so each one is
%  written to produce both good print output and clean HTML:
%
%    \entry  -> <h3><a>organisation</a><em>dates</em></h3>
%               <p><strong>role</strong></p>
%    \project / \projectplain -> a run-in bold title starting a paragraph
%
%  The \hfill in \entry right-aligns the dates in the PDF; pandoc drops
%  it, and the stylesheet right-aligns the <em> instead.
% ---------------------------------------------------------------------

% \entry{url}{organisation}{dates}{role}
\newcommand{\entry}[4]{%
  \subsection*{\href{#1}{#2}\hfill\emph{#3}}%
  \textbf{#4}\par
}

% \project{url}{Title:} -- run-in title that links out.
% The newline after the call supplies the following space; do not add a
% trailing `\ ', which pandoc would turn into a non-breaking space.
\newcommand{\project}[2]{\textbf{\href{#1}{#2}}}

% \projectplain{Title:} -- run-in title with no link
\newcommand{\projectplain}[1]{\textbf{#1}}

\begin{document}
\maketitle

\begin{center}
York, United Kingdom \textbullet{}
+44 7941 079483 \textbullet{}
\href{mailto:jonathanwoollettlight@proton.me}{jonathanwoollettlight@proton.me}
\\
\href{https://github.com/JonathanWoollett-Light}{github.com/JonathanWoollett-Light}
\textbullet{}
\href{https://www.linkedin.com/in/jonathan-wl/}{linkedin.com/in/jonathan-wl}
\textbullet{}
\href{https://www.youtube.com/@alittlemorethananintroduct1604}{YouTube}
\end{center}

I am a PhD student in Physics at the University of York with the research topic
of `Neuromorphic Computing for Resource-Efficient Autonomous Systems'. This
entails investigation of a wide range of AI/ML/RL algorithms and their
respective hardware accelerators with particular focus and consideration of
neurmorphic hardware and algorithms, and how they compare to more traditional
ANNs. I have a specific interest in low level systems and tooling. I am a
frequent Open-Source contributor having contributed to ziglang, rustlang and
popular Rust crates. I have a small educational YouTube channel where I very
infrequently post videos about technical topics I find interesting
(\href{https://www.youtube.com/@alittlemorethananintroduct1604}{A little more
than an introduction to -- YouTube}).

\section{Experience}

\entry{https://www.york.ac.uk/}{York University}{April 2025 -- Current}{PhD in Physics}
I am currently undertaking a PhD in Physics at the University of York with the
research topic of `Neuromorphic Computing for Resource-Efficient Autonomous
Systems'.

\entry{https://aws.amazon.com/}{Amazon Web Services}{April 2022 -- June 2024}{\href{https://maps.app.goo.gl/Drok6UA7W5empdoY6}{Cambrigde} -- Software development engineer}
I worked on Firecracker, a memory efficient, fast booting VMM that allows AWS
lambda to be profitable.

Firecracker supports snapshotting a VM; a VM can be saved and restored much
faster than booting a new one and much faster than any other VMM. I spent time
working on the safety and configurability of this feature, this involved
evaluating whether a snapshot could be safely restored on a given CPU
considering its specified features and security mitigations, and those of the
target. This involved analysis of the performance and security effects of CPU
features on both x86-64 and aarch64, and ensuring we had an implementation that
aligned with the specifications.

Firecracker supports a feature called CPU templates that allows a user
configurability of most of how the CPU is presented to the guest OS. This work
required understanding the Intel x86-64, AMD x86-64 and Arm aarch64
specifications. Development involved thorough discussion with customers and
within our team both focused on usability, security and robustness. I rewrote
how Firecracker handles static CPU templates to support runtime CPU templates
and wrote the foundational implementation for runtime CPU templates. I also
implemented thorough compatibility checks for x86-64.

I was the owner of an initiative to performance test Firecracker in
production-like scenarios, through this we discovered and fixed significant
performance regressions that may have costed AWS Lambda millions of dollars if
they went live. This work involved coordinating across teams and diving deep
into the deployment and management of resources to setup the required
environment and execute the useful tests.

In this role I and another team member adopted responsibility for maintenance of
the \href{https://github.com/rust-vmm}{rust-vmm} open source organization that
maintained components used by Firecracker and other VMMs such as
\href{https://github.com/cloud-hypervisor/cloud-hypervisor}{cloud-hypervisor}.

\entry{https://www.renishaw.com/}{Renishaw}{July 2018 -- September 2018}{\href{https://maps.app.goo.gl/d1S4cdULwPKQcSf98}{Wotton-under-Edge} -- Software placement}
I worked on building C++ components forming an internal standard library. This
involved extensive testing and performance consideration.

\section{Education}

\entry{https://www.swansea.ac.uk/}{Swansea University}{2017 -- 2021}{MEng Computing \& Bachelors 2.1}
I competed with the university Taekwondo society, at times internationally.

\entry{https://www.renishaw.com/}{Renishaw}{A few weeks in 2016}{\href{https://maps.app.goo.gl/d1S4cdULwPKQcSf98}{Wotton-under-Edge} -- Work experience}
My work experience focused on development with FPGAs and VHDL.

\entry{https://www.bhasvic.ac.uk/}{Brighton and Hove Sixth Form College}{2015 -- 2017}{A levels -- Mathematics: B, Modern History: B, Computer Science: C}
I did an Extended Project Qualification discussing the question `Is the
development and completion of an AGI (Artificial General Intelligence) likely
within the next 80 years?'.

\entry{https://en.wikipedia.org/wiki/St_Andrew\%27s_High_School,_Worthing}{St Andrew's High School}{2011 -- 2015}{GCSEs 11 A*--C including mathematics and english}
I founded the computing club, the debating club and was elected deputy head boy
by students.

\section{Portfolio}

I am an expert in Rust, I am experienced in C++, C and Python, and I am familiar
with Zig, Ada, aarch64 asm, RISC-V asm, JavaScript and CSS3+HTML5.

\project{https://jonathanwoollett-light.github.io/formal/}{My programming language:}
A verifying compiler for bare-metal RISC-V, written in Rust. You write a small
Python-like language, and the compiler accepts your program only if it can prove
(by symbolically executing the actual machine code across every hardware-thread
interleaving and every admissible type assignment) that no assertion can fail
and no memory access is ever out of bounds, even for intentionally racy,
lock-free code. Its type inference is complete, or `infallible': a variable left
unspecified is accepted exactly when some typing makes the whole program verify,
because inference and verification are the same search. The proof's by-products
then drive code generation: it infers the memory layout and performs dead-data
elimination, deleting interior bytes of live data it proved are never touched at
runtime and rewriting the address arithmetic to match. It verifies the
instructions themselves rather than high-level source, shrinking the trusted
base to a small layout-only lowering plus the assembler and linker; conceptually
it generalises Zig's concrete compile-time evaluation to full symbolic
exploration and turns Rust's conservative data-race prohibition into exact,
per-assertion race verification, with parallel and distributed (MPI) verifiers
checked against a trusted single-threaded oracle.
(\href{https://github.com/JonathanWoollett-Light/formal}{source})

\projectplain{CV website:}
A simple website serving the latest version of my curriculum vitae. Written in
LaTeX and rendered to a single self-contained page: a near-perfect Google
Lighthouse score, light and dark mode support, less than 100kb in size, and
neatly prints to PDF.

\projectplain{Hearts of Iorn 4 modding:}
I'm the owner and contributor to multiple Hearts of Iron 4 mods on the Steam
Workshop
(\href{https://steamcommunity.com/sharedfiles/filedetails/?id=3342313594}{OWB --
The Think Tank} and
\href{https://steamcommunity.com/sharedfiles/filedetails/?id=2777392649\&searchtext=}{Millennium
Dawn: A Modern Day Mod}). Working with mods and the often non-technical people
making them requires a uniquely scrappy approach to development and a lack of
hangups on formal processes (being able to work within what might politely be
described as a \emph{mess}), certainly a mindset shift from my traditional
systems development.

\projectplain{Dating app:}
A dating app that learns user specific traits and preferences to match pairs.
The server functions as an efficient monolith capable of handling all requests
on a single Rust \href{https://github.com/tokio-rs/axum}{Axum} server with
in-memory stores and a local SQL store, achieving near 100\% system utilisation
capable of handling millions of users; supporting all the functionality of a
typical tinder-like dating app.

\project{https://github.com/ci-metrics/example}{CI Metrics:}
A website to track performance metrics (continuous benchmarking) and integrate
with GitHub. A Rust \href{https://github.com/tokio-rs/axum}{Axum} backend with a
MongoDB database. Development focussed on experimentation with algorithms for
anomaly detection to produce a solution that accomplishes what
\href{https://aws.amazon.com/lookout-for-metrics/}{AWS Lookout for Metrics}
does, but at a much lower cost and complexity. This was motivated by my
experience working on
\href{https://github.com/firecracker-microvm/firecracker}{Firecracker} at AWS
where regressions that could have cost millions in production were almost
missed, and certainly some were.

\projectplain{Live migration web server:}
A Rust server framework where a new server process can succeed a local old
process in milliseconds with shared memory (or remotely in seconds over TCP)
without dropping client streams. This can be used as a database server with
in-memory data storage akin to Redis where queries are compiled as new
functions.

\project{https://github.com/JonathanWoollett-Light/rust-ad}{RUST-AD:}
Auto-differentiation library with code transformation for both forward and
reverse auto-diff written in Rust.

\projectplain{Mathematics Site:}
A mathematics website for automated marking. Built with a Rust
\href{https://github.com/actix/actix-web}{Actix} backend using a PyTorch
semantic segmentation model, deployed with AWS, with a MongoDB database.

\project{https://github.com/JonathanWoollett-Light/GLSL-BLAS}{GLSL BLAS:}
GLSL code for non-complex BLAS operations with Vulkan. My final year masters
project. Using modern C++ compile-time evaluation to validate shaders.

\project{https://jonathanwoollett-light.github.io/BLAS/}{BLAS website:}
A website documenting BLAS operations. While working on GLSL BLAS I found
\href{https://www.netlib.org/blas/}{the older website} used to document BLAS to
be a little uncomfortable, so I made this. I've also added some Python
\href{https://www.manim.community/}{Manim} animations to illustrate the
operations.

\project{https://github.com/JonathanWoollett-Light/cogent}{Cogent:}
A simple library Rust library using \href{https://arrayfire.com/}{ArrayFire} for
training simple neural networks; like a small Keras. A component of my final
year bachelors project.

\projectplain{Genetic Draughts AI:}
My first ever venture into machine learning and optimisation. A simple tree C++
search program to play draughts using a genetic algorithm to judge board
positions. I still remember the moment it beat me.

\end{document}
